% ============================================================
% SECTION 02
% IDENTITY PRINCIPLES AND PROPOSITIONS
% Chapter 06 — Synthesis of Identity
% Book I: Ontology of Viability
% ============================================================

\section{Identity Principles and Propositions}
\label{sec:identity-principles-and-propositions}

% ------------------------------------------------------------
% 2.1 — THE STATUS OF PRINCIPLES AND PROPOSITIONS
% ------------------------------------------------------------

\subsection{The Status of Principles and Propositions}

The integrated account developed in the preceding section
provides the basis for distinguishing the different logical
and ontological statuses of the claims concerning identity.
Not every statement established during the preceding
investigations has the same degree of generality, necessity,
or evidential support. Some statements express foundational
commitments of the present ontological framework, whereas
others formulate distinctions, restrictions, conditional
relations, or questions whose validity remains dependent
upon further determination.

The distinction between principles and propositions is
therefore essential. A principle establishes a governing
commitment that constrains the interpretation of identity
throughout the inquiry. A proposition articulates a more
specific claim concerning identity or its relation to
another ontological phenomenon. Propositions may follow
from principles, depend upon additional premises, or remain
valid only under specified conditions.

This distinction must not be confused with the difference
between truth and hypothesis. A proposition may express
an established ontological distinction, a conditional
claim, or a hypothesis requiring further examination.
Likewise, a principle must not be treated as established
merely because it has been given foundational terminology.
Its status depends upon its justification within the
ontological argument.

The purpose of this section is to formulate the governing
principles of the present account, identify the propositions
that follow from the preceding investigations, and clarify
the limits within which these claims may be asserted.

% ------------------------------------------------------------
% 2.2 — FOUNDATIONAL PRINCIPLES
% ------------------------------------------------------------

\subsection{Foundational Principles}

The following principles express the principal ontological
commitments established through the inquiry into identity.
They are intended to govern subsequent analysis without
prematurely resolving questions that remain open.

\subsubsection{Principle of Ontological Inseparability}

Identity and existence are ontologically inseparable.
An existent is not first established as an identity-free
entity to which identity is subsequently added. Nor is
identity an independently existing determination that
can be detached from the existent whose identity is at
issue.

This inseparability does not eliminate their analytical
distinction. The question of existence and the question
of identity remain distinguishable inquiries, although
neither establishes an ontological priority over the other.

\subsubsection{Principle of Analytical Distinction}

The analytical distinction between identity and its
associated ontological phenomena must be preserved.
Determination, invariance, unity, singularity, composition,
expression, manifestation, persistence, and continuity
may be relevant to the examination of identity, but none
may be substituted for identity without an independent
ontological justification.

This principle prohibits conceptual reduction while
allowing the investigation of relations between identity
and other dimensions of an existent.

\subsubsection{Principle of Non-Reduction}

Identity cannot be exhaustively defined through a single
feature, relation, structure, invariant, manifestation,
or continuity condition.

The relevance of a phenomenon to identity does not
establish its equivalence with identity. Any proposed
reduction must therefore demonstrate that the phenomenon
in question is ontologically sufficient for the account
being advanced, rather than relying upon terminological
identification.

\subsubsection{Principle of Subject and Respect Specification}

Every ontological proposition concerning identity must
be evaluated with respect to the existent to which it
applies and the particular respect under which the claim
is made.

A proposition established concerning material constitution,
structural organization, relational determination, or
manifestation cannot automatically be generalized to
other respects. Its scope must remain within the domain
and conditions under which its ontological relevance
has been established.

\subsubsection{Principle of Ontological Independence from Recognition}

The identity of an existent does not depend upon its
recognition, observation, description, or representation
by an observer.

Epistemic access may affect what can be established about
identity, but it does not constitute the ontological
conditions of identity itself. The absence of recognition
or identification must not be treated as evidence of
ontological absence or non-identity without further
argument.

\subsubsection{Principle of Non-Presumption}

The number, unity, or plurality of identities associated
with an existent must not be assumed in advance.

Neither unity nor plurality may be assigned universal
ontological priority without sufficient justification.
The question must remain open wherever the relevant
conditions have not been established.

% ------------------------------------------------------------
% 2.3 — PRINCIPLES GOVERNING INVARIANCE AND CHANGE
% ------------------------------------------------------------

\subsection{Principles Governing Invariance and Change}

The inquiry into invariance and continuity established
that the relation between identity and change cannot
be governed by a universal assumption that every
transformation either preserves or destroys identity.

\subsubsection{Principle of the Relativity of Invariance}

Every claim of invariance must specify the transformation,
domain, and respect relative to which something is said
to remain invariant.

An invariant identified under one transformation does
not necessarily remain invariant under another. Nor
does formal invariance independently establish ontological
invariance.

\subsubsection{Principle of the Non-Equivalence of Identity and Invariance}

The preservation of an invariant does not independently
establish the persistence of identity. Likewise, the
absence of an identified invariant does not independently
establish non-identity.

An invariant may become relevant to identity preservation
only when its ontological significance for the existent
under examination has been established.

\subsubsection{Principle of the Non-Equivalence of Change and Non-Identity}

A difference between two conditions, stages, or
manifestations does not independently establish numerical
non-identity.

Qualitative change, material alteration, structural
modification, relational reconfiguration, and variation
in manifestation must be distinguished from the
ontological termination or replacement of an identity.

\subsubsection{Principle of the Non-Equivalence of Continuity and Identity}

Temporal succession or continuity of manifestations
does not independently establish the persistence of
identity.

Continuity must be examined through its specific subject,
domain, and ontological respect. The existence of a
continuous sequence does not by itself demonstrate that
the sequence belongs to one numerically identical
existent or identity.

% ------------------------------------------------------------
% 2.4 — PRINCIPLES GOVERNING UNITY AND PLURALITY
% ------------------------------------------------------------

\subsection{Principles Governing Unity and Plurality}

The investigation of unity and plurality imposes further
restrictions upon the formulation of identity propositions.
In particular, the unity of an existent cannot be inferred
from a single form of unity, nor can plurality be inferred
from multiplicity in an unrelated ontological respect.

\subsubsection{Principle of the Distinction Between Forms of Unity}

Unity of identity, unity of source, and unity of aggregate
are distinct ontological questions.

The establishment of one form of unity does not
automatically establish another. Their relations require
independent examination.

\subsubsection{Principle of the Non-Equivalence of Unity and Singularity}

Unity must not be identified with singularity.

A unified existent or condition may involve multiplicity
without being reducible to an aggregate. Conversely,
singularity cannot be inferred merely from the presence
of unity.

\subsubsection{Principle of the Non-Equivalence of Multiplicity and Identity Plurality}

The plurality of components, relations, determinations,
expressions, or manifestations does not independently
establish a plurality of identities.

The ontological relation between multiplicity and identity
plurality must be established rather than assumed.

\subsubsection{Principle of the Open Status of Identity Plurality}

Neither the thesis that every existent possesses exactly
one identity nor the thesis that an existent may possess
multiple identities has been established as a universal
principle.

The number of identities associated with an existent
must therefore remain subject to further ontological
determination.

% ------------------------------------------------------------
% 2.5 — PRINCIPLES GOVERNING EXPRESSION AND MANIFESTATION
% ------------------------------------------------------------

\subsection{Principles Governing Expression and Manifestation}

The distinction between identity, expression, and
manifestation constrains the interpretation of how an
existent is expressed or manifested without treating
these phenomena as substitutes for its identity.

\subsubsection{Principle of the Non-Equivalence of Identity and Expression}

Expression is not identical with the identity of that
which is expressed.

The presence, absence, variation, or plurality of
expressions does not independently establish the
preservation, termination, or plurality of identity.

\subsubsection{Principle of the Non-Equivalence of Expression and Manifestation}

Expression and manifestation are distinct ontological
problems. Their relation must be investigated without
assuming that every expression is a manifestation or
that every manifestation exhausts the phenomenon of
expression.

\subsubsection{Principle of the Non-Equivalence of Manifestation and Identity}

A particular manifestation cannot be identified with
the identity of the existent associated with it.

The termination or alteration of a manifestation does
not independently establish the termination or alteration
of identity. Such a conclusion requires an account of
the ontological relation between identity and the
manifestational conditions in question.

\subsubsection{Principle of the Distinction Between Manifestation and Capacity to Manifest}

The continuity of actual manifestations must be
distinguished from the continuity of the capacity to
manifest.

The absence of an actual manifestation does not
independently establish the absence of such a capacity,
and the existence of a capacity cannot be inferred
solely from the absence of manifestation.

% ------------------------------------------------------------
% 2.6 — PROPOSITIONS DERIVED FROM THE INQUIRY
% ------------------------------------------------------------

\subsection{Propositions Derived from the Inquiry}

The principles formulated above support a series of
propositions concerning the relations between identity
and the phenomena examined in the preceding chapters.
These propositions express the principal consequences
of the present framework. Their scope remains restricted
by the ontological conditions and distinctions already
established.

\subsubsection{Proposition 1: Relevance Does Not Entail Equivalence}

If a phenomenon is relevant to the constitution,
preservation, expression, or manifestation of an identity,
it does not follow that the phenomenon is identical
with that identity.

\subsubsection{Proposition 2: Preservation of Features Does Not Establish Identity}

The preservation of one or more features of an existent
does not, by itself, establish the preservation of its
numerical identity.

\subsubsection{Proposition 3: Absence of an Invariant Does Not Establish Non-Identity}

The failure to identify an invariant under specified
conditions does not, by itself, establish that identity
has been terminated or replaced.

\subsubsection{Proposition 4: Difference Does Not Establish Numerical Non-Identity}

A qualitative difference between two states or
manifestations does not, without additional ontological
conditions, establish that they belong to different
numerical identities.

\subsubsection{Proposition 5: Continuity Does Not Independently Establish Identity Persistence}

The continuity of a temporal sequence, causal relation,
or manifestation does not independently establish the
persistence of a single identity throughout that sequence.

\subsubsection{Proposition 6: Interruption Does Not Independently Establish Identity Termination}

An interruption in manifestation, activity, or an
established relation does not independently establish
the termination of the identity associated with the
interrupted condition.

\subsubsection{Proposition 7: Multiplicity Does Not Establish Identity Plurality}

The multiplicity of components, relations, expressions,
or manifestations does not independently establish
that an existent possesses multiple identities.

\subsubsection{Proposition 8: Unity Does Not Establish Identity Singularity}

The unity of an existent, source, or aggregate does
not independently establish that only one numerical
identity is associated with it.

\subsubsection{Proposition 9: Ontological Claims Require Restricted Scope}

A proposition concerning identity established under
a particular domain, transformation, or respect cannot
be generalized beyond those conditions without further
justification.

\subsubsection{Proposition 10: Epistemic Insufficiency Does Not Establish Ontological Absence}

The inability to recognize, describe, or determine an
identity does not independently establish that the
identity is absent or nonexistent.

% ------------------------------------------------------------
% 2.7 — CONDITIONAL PROPOSITIONS
% ------------------------------------------------------------

\subsection{Conditional Propositions}

Some claims concerning identity cannot be asserted
without specifying additional ontological conditions.
They are therefore formulated conditionally rather than
as universal principles.

First, an invariant may serve as a condition of identity
preservation if its ontological relevance to the existent
and transformation under examination has been established.
Its formal preservation alone is insufficient.

Second, a material, structural, relational, or causal
condition may contribute to identity preservation if
the relevant ontological dependence has been demonstrated.
No such condition should be treated as universally
necessary or sufficient without establishing its scope.

Third, continuity of manifestation may support a claim
of identity continuity if the ontological relation between
the successive manifestations and the identity in
question has been independently established.

Fourth, the unity of an existent may support a claim
concerning the unity of identity if the relevant form
of unity is shown to be ontologically constitutive of
identity rather than merely associated with its source,
composition, or manifestation.

Fifth, the plurality of expressions or manifestations
may be relevant to the question of identity plurality
if an ontological relation between those expressions
or manifestations and distinct identities can be
established.

These conditional propositions do not establish the
truth of the conditions they specify. They identify
the additional requirements under which more specific
claims may become defensible.

% ------------------------------------------------------------
% 2.8 — METHODOLOGICAL RESTRICTIONS
% ------------------------------------------------------------

\subsection{Methodological Restrictions}

The principles and propositions of identity impose
methodological restrictions upon subsequent ontological
inquiry.

First, conceptual distinctions must not be eliminated
merely because two phenomena frequently occur together.
Co-occurrence does not establish ontological equivalence.

Second, a relation established within one domain must
not be generalized to all domains without sufficient
justification. The universality of a condition is itself
a claim requiring argument.

Third, the absence of an established condition must
not be converted into proof of its ontological absence.
The failure of an inquiry to establish continuity,
invariance, or identity plurality may reflect the limits
of available determination rather than a definitive
ontological result.

Fourth, necessary and sufficient conditions must be
distinguished. A condition shown to be necessary under
specified circumstances is not thereby sufficient, and
a sufficient condition within a restricted domain is
not thereby universally necessary.

Fifth, formal models and scientific descriptions may
constrain or clarify ontological propositions, but their
formal validity does not independently establish the
universal ontological status of their conclusions.

Finally, unresolved questions must remain explicitly
unresolved. A theoretical framework should not obtain
apparent completeness by converting analytical
insufficiency into an unsupported ontological assertion.

% ------------------------------------------------------------
% 2.9 — THE STATUS OF THE FORMULATED FRAMEWORK
% ------------------------------------------------------------

\subsection{The Status of the Formulated Framework}

The principles and propositions presented in this
section constitute a structured account of the claims
supported by the preceding inquiry. They do not
represent a complete axiomatic system from which every
fact concerning identity can already be derived.

The principles establish restrictions and governing
commitments for the ontological analysis of identity.
The propositions articulate consequences of those
commitments. The conditional propositions specify
relations that may hold under additional conditions.
The methodological restrictions determine how these
claims may be applied and generalized.

Their formulation must therefore not be mistaken for
a proof that every ontological problem of identity has
been resolved. The framework is sufficiently determinate
to constrain subsequent inquiry, while remaining open
to revision where further ontological investigation
discloses inadequacies, counterexamples, or additional
distinctions.

The distinction between established claims and open
questions is not a weakness of the framework. It is
a requirement of ontological rigor: the scope of a
claim must remain proportionate to the justification
available for it.

% ------------------------------------------------------------
% 2.10 — SCOPE OF THE PRESENT SECTION
% ------------------------------------------------------------

\subsection{Scope of the Present Section}

This section has formulated the principal principles
and propositions governing the present account of
identity. It has distinguished foundational commitments
from derived propositions and conditional claims, while
establishing restrictions upon their scope and
application.

It has not attempted to settle every question concerning
the ontological constitution, plurality, or preservation
of identity. Nor has it converted the distinctions
developed in the preceding chapters into a closed
system that excludes further ontological investigation.

The next section will identify the principal questions
that remain open. Its purpose will be to distinguish
problems that require further ontological determination
from limitations arising from the present framework
or from the conditions under which particular claims
can be established.

% ============================================================
% END OF SECTION 02
% ============================================================